\documentclass[11pt,a4paper]{article}

\usepackage{amsmath, amssymb, amsthm}
\usepackage{physics}
\usepackage{bm}
\usepackage{siunitx}
\usepackage{booktabs}
\usepackage{enumitem}
\usepackage[margin=1in]{geometry}
\usepackage{hyperref}

\title{Technical Note: Data Inversion Theory for the DMA-APM-CPC\\
Tandem Aerosol Measurement System}
\author{N.~Moteki}
\date{Last updated: 2026-09-24}

\begin{document}
\maketitle

\begin{abstract}
This document provides a self-contained description of the mathematical
theory and numerical algorithms implemented in the DMA-APM mass distribution
inversion tool.  It covers the formulation of the forward problem (1D and 2D
integral models), the physical justification for the 1D approximation, the
original RK4-based APM transfer function simulation, the Chahine--Twomey
inversion algorithm with Markowski smoothing, and the Poisson-statistics-based
convergence criterion.  It also describes the use of the standalone transfer
function for planning APM set-points (rotation speed and voltage) for
spherical particles and multiply charged sphere clusters.
\end{abstract}

\tableofcontents
\newpage

%% ====================================================================
\section{Objective and problem statement}
\label{sec:objective}
%% ====================================================================

The objective of this analysis is to estimate the highly resolved mass
distribution of aerosol particles that have been pre-classified by a specific
electrical mobility diameter using a DMA-APM-CPC tandem measurement system.
Specifically, we seek the mass distribution function at the DMA outlet:
\begin{equation}
  f(m) = \frac{\dd N}{\dd m},
  \label{eq:mass-distribution}
\end{equation}
where $N$ is the particle number concentration and $m$ is the particle mass.
The observed CPC concentration $n(V)$ at each APM applied voltage $V$ is
related to $f(m)$ through a Fredholm integral equation of the first kind, and
recovering $f(m)$ from $n(V)$ constitutes an ill-posed inverse problem.

%% ====================================================================
\section{Formulation of the forward problem}
\label{sec:forward-problem}
%% ====================================================================

%% --------------------------------------------------------------------
\subsection{The rigorous 2D-integral model}
\label{sec:2d-model}
%% --------------------------------------------------------------------

The expected particle number concentration $n(V)$ observed downstream of the
APM is governed by a double integral over both mass $m$ and electrical
mobility $Z_{\mathrm{p}}$:
\begin{equation}
  n(V) = \int_{0}^{\infty}\!\int_{0}^{\infty}
    \Omega_{\mathrm{APM}}(m, V, Z_{\mathrm{p}}) \,
    f_{\mathrm{in}}(m, Z_{\mathrm{p}}) \,
    \Omega_{\mathrm{DMA}}(Z_{\mathrm{p}}) \,
    \dd Z_{\mathrm{p}} \, \dd m,
  \label{eq:2d-forward}
\end{equation}
where $f_{\mathrm{in}}(m, Z_{\mathrm{p}})$ is the intrinsic 2D
mass-mobility distribution of the raw aerosol entering the DMA,
$\Omega_{\mathrm{DMA}}(Z_{\mathrm{p}})$ is the DMA transfer function, and
$\Omega_{\mathrm{APM}}(m, V, Z_{\mathrm{p}})$ is the APM transfer function.

Assuming that $f_{\mathrm{in}}$ is approximately constant with respect to
$Z_{\mathrm{p}}$ over the narrow DMA transmission window, we define the
target distribution $f(m)$ as the mass distribution of the aerosol exiting
the DMA:
\begin{equation}
  f(m) = f_{\mathrm{in}}(m, Z_{\mathrm{p}}^{*})
    \int_{0}^{\infty} \Omega_{\mathrm{DMA}}(Z_{\mathrm{p}}) \,
    \dd Z_{\mathrm{p}}.
  \label{eq:target-distribution}
\end{equation}

Using this definition, the forward problem reduces to:
\begin{equation}
  n(V) = \int_{0}^{\infty}
    K_{\mathrm{eff}}^{\mathrm{2D}}(V, m) \, f(m) \, \dd m,
  \label{eq:2d-reduced}
\end{equation}
where the effective 2D kernel is:
\begin{equation}
  K_{\mathrm{eff}}^{\mathrm{2D}}(V, m) =
    \frac{\displaystyle\int_{0}^{\infty}
      \Omega_{\mathrm{APM}}(m, V, Z_{\mathrm{p}}) \,
      \Omega_{\mathrm{DMA}}(Z_{\mathrm{p}}) \,
      \dd Z_{\mathrm{p}}}
    {\displaystyle\int_{0}^{\infty}
      \Omega_{\mathrm{DMA}}(Z_{\mathrm{p}}) \,
      \dd Z_{\mathrm{p}}}.
  \label{eq:2d-kernel}
\end{equation}

In the implementation, the DMA transfer function is modelled as a triangular
function centred at the target mobility $Z_{\mathrm{p}}^{*}$ with half-width
$\beta Z_{\mathrm{p}}^{*}$, where
$\beta = Q_{\mathrm{sample}}/Q_{\mathrm{sheath}}$ is the DMA flow ratio.
The integral in \eqref{eq:2d-kernel} is evaluated numerically using the
rectangle rule with a configurable number of quadrature points.

%% --------------------------------------------------------------------
\subsection{The 1D-integral approximation}
\label{sec:1d-model}
%% --------------------------------------------------------------------

By approximating the DMA output as a Dirac delta function centred at the
target mobility $Z_{\mathrm{p}}^{*}$, the effective kernel simplifies to:
\begin{equation}
  K^{\mathrm{1D}}(V, m) =
    \Omega_{\mathrm{APM}}(m, V, Z_{\mathrm{p}}^{*}),
  \label{eq:1d-kernel}
\end{equation}
and the forward problem becomes:
\begin{equation}
  n(V) \approx \int_{0}^{\infty}
    K^{\mathrm{1D}}(V, m) \, f(m) \, \dd m.
  \label{eq:1d-forward}
\end{equation}

In both the 1D and 2D models, the inversion algorithm solves for the same
target: $f(m)$, the mass distribution of the particles at the DMA outlet.

%% ====================================================================
\section{Physical justification: why the 1D model is practically sufficient}
\label{sec:1d-justification}
%% ====================================================================

In practice, the mass distributions obtained from the 1D and 2D models are
virtually identical, even when the DMA resolution is relatively broad
(e.g., $\beta = 0.2$).  This equivalence arises from the fundamental
classification principle of the APM.  The central classified mass is:
\begin{equation}
  m_{\mathrm{c}} =
    \frac{e V}
         {r_{\mathrm{c}}^{2} \, \omega^{2} \, \ln(r_{2}/r_{1})},
  \label{eq:central-mass}
\end{equation}
where $e$ is the elementary charge,
$r_{\mathrm{c}} = (r_{1} + r_{2})/2$ is the gap centre radius,
$\omega$ is the angular velocity, and $r_{1}$, $r_{2}$ are the inner and
outer cylinder radii.  This expression is completely independent of the
particle's electrical mobility $Z_{\mathrm{p}}$.  A variation in
$Z_{\mathrm{p}}$ (due to the DMA's finite transmission width) only alters the
particle's radial drift velocity inside the APM, slightly affecting its
transit time.  Consequently, $Z_{\mathrm{p}}$ variations marginally broaden
the width of $\Omega_{\mathrm{APM}}$ but do not shift its central peak.
Because the 2D convolution averages these minor, symmetric width changes, the
resulting $K_{\mathrm{eff}}^{\mathrm{2D}}$ is nearly indistinguishable from
$K^{\mathrm{1D}}$.

The 1D approximation therefore provides mathematically and physically robust
results at approximately $1/N_{Z_{\mathrm{p}}}$ of the computational cost of
the 2D model, where $N_{Z_{\mathrm{p}}}$ is the number of $Z_{\mathrm{p}}$
quadrature points.

%% ====================================================================
\section{APM transfer function: particle trajectory simulation}
\label{sec:apm-transfer}
%% ====================================================================

%% --------------------------------------------------------------------
\subsection{Equation of motion}
\label{sec:equation-of-motion}
%% --------------------------------------------------------------------

To compute $\Omega_{\mathrm{APM}}(m, V, Z_{\mathrm{p}})$, this tool employs
an original numerical simulation method that directly tracks particle
trajectories within the APM gap using the 4th-order Runge--Kutta (RK4) method.

Inside the APM annular gap (inner radius $r_{1}$, outer radius $r_{2}$,
centre radius $r_{\mathrm{c}} = (r_{1} + r_{2})/2$, half-width
$\delta = (r_{2} - r_{1})/2$), a particle of mass $m$ and electrical
mobility $Z_{\mathrm{p}}$ experiences centrifugal and electrostatic forces.
Under the assumption of a fully developed parabolic laminar flow profile, the
ratio of the radial to axial velocity components is
(Ehara et al., 1996):
\begin{equation}
  \frac{\dd r}{\dd z}
  = \frac{v_{\mathrm{r}}(r)}{v_{\mathrm{z}}(r)}
  = \frac{\dfrac{Z_{\mathrm{p}}}{e}
      \left( m \omega^{2} r
             - \dfrac{eV}{r \ln(r_{2}/r_{1})} \right)}
    {\dfrac{3Q}{4\pi \delta r_{\mathrm{c}}}
      \left[ 1 - \left(\dfrac{r - r_{\mathrm{c}}}{\delta}\right)^{2}
      \right]},
  \label{eq:motion-raw}
\end{equation}
where $Q$ is the volumetric aerosol flow rate.  This can be rearranged to:
\begin{equation}
  \frac{\dd r}{\dd z}
  = \frac{8\pi \delta r_{\mathrm{c}}}{3Q}
    \frac{Z_{\mathrm{p}}}{e} \cdot
    \frac{m \omega^{2} r
          - \dfrac{eV}{r \ln(r_{2}/r_{1})}}
         {1 - \left(\dfrac{r - r_{\mathrm{c}}}{\delta}\right)^{2}}.
  \label{eq:motion-rearranged}
\end{equation}

%% --------------------------------------------------------------------
\subsection{RK4 integration and transmission efficiency}
\label{sec:rk4}
%% --------------------------------------------------------------------

The trajectory equation~\eqref{eq:motion-rearranged} is integrated along the
axial coordinate $z$ from $z = 0$ to $z = L$ (the electrode length) using
the classical 4th-order Runge--Kutta method with a fixed step size $\Delta z$:
\begin{align}
  k_{1} &= \Delta z \cdot g(r_{n}),
    \label{eq:rk4-k1} \\
  k_{2} &= \Delta z \cdot g(r_{n} + k_{1}/2),
    \label{eq:rk4-k2} \\
  k_{3} &= \Delta z \cdot g(r_{n} + k_{2}/2),
    \label{eq:rk4-k3} \\
  k_{4} &= \Delta z \cdot g(r_{n} + k_{3}),
    \label{eq:rk4-k4} \\
  r_{n+1} &= r_{n} + \frac{1}{6}(k_{1} + 2k_{2} + 2k_{3} + k_{4}),
    \label{eq:rk4-update}
\end{align}
where $g(r) \equiv \dd r / \dd z$ as defined in
\eqref{eq:motion-rearranged}.

A total of $N_{r_{0}}$ particles are launched from uniformly spaced initial
radial positions $r_{0} \in [r_{1}, r_{2}]$.  A particle is considered lost
if it hits the inner or outer electrode wall at any step.  The transmission
efficiency is computed as the flux-weighted fraction of surviving particles:
\begin{equation}
  \Omega_{\mathrm{APM}}
  = \frac{\displaystyle\sum_{r_{0} \in \text{survived}} w(r_{0})}
         {\displaystyle\sum_{r_{0} \in \text{all}} w(r_{0})},
  \label{eq:transmission}
\end{equation}
where the weight $w(r_{0})$ accounts for the parabolic flow profile:
\begin{equation}
  w(r_{0}) = \frac{3}{2}
    \left[ 1 - \left(\frac{r_{0} - r_{\mathrm{c}}}{\delta}\right)^{2}
    \right].
  \label{eq:flux-weight}
\end{equation}

%% --------------------------------------------------------------------
\subsection{Cunningham slip correction}
\label{sec:cunningham}
%% --------------------------------------------------------------------

The electrical mobility $Z_{\mathrm{p}}$ is related to the mobility diameter
$D_{\mathrm{mob}}$ via:
\begin{equation}
  Z_{\mathrm{p}} =
    \frac{e \, C_{\mathrm{c}}(D_{\mathrm{mob}})}
         {3 \pi \eta D_{\mathrm{mob}}},
  \label{eq:mobility}
\end{equation}
where $\eta$ is the dynamic viscosity of air and $C_{\mathrm{c}}$ is the
Cunningham slip correction factor (Hinds, 1999):
\begin{equation}
  C_{\mathrm{c}} = 1 + \frac{1}{Pd}
    \left[ 15.60 + 7.00 \exp(-0.059 \, Pd) \right],
  \label{eq:cunningham}
\end{equation}
with $P$ the atmospheric pressure in \si{\kilo\pascal} and $d$ the diameter
in \si{\micro\metre}.

%% --------------------------------------------------------------------
\subsection{Kernel matrix discretisation}
\label{sec:kernel-discretisation}
%% --------------------------------------------------------------------

For the discrete inversion, the integral in \eqref{eq:1d-forward} is
approximated by the rectangle rule over $J$ mass bins:
\begin{equation}
  n(V_{i}) \approx \sum_{j=1}^{J} K_{i,j} \, f(m_{j}),
  \label{eq:discrete-forward}
\end{equation}
where
$K_{i,j} = \Omega_{\mathrm{APM}}(m_{j}, V_{i}, Z_{\mathrm{p}}^{*})
\cdot \Delta m$
for the 1D model, $m_{j}$ are uniformly spaced mass grid points, and
$\Delta m = (m_{\mathrm{max}} - m_{\mathrm{min}})/(J - 1)$.
The kernel matrix $\mathbf{K}$ has shape $I \times J$, where $I$ is the
number of voltage bins.

%% ====================================================================
\section{Inversion algorithm: Chahine--Twomey method with internal smoothing}
\label{sec:inversion}
%% ====================================================================

%% --------------------------------------------------------------------
\subsection{Update equation}
\label{sec:update-equation}
%% --------------------------------------------------------------------

To solve the discrete ill-posed problem
$\mathbf{n} = \mathbf{K} \mathbf{f}$, we use the non-negative
Chahine--Twomey iterative algorithm (Twomey, 1975).  The update rule at
iteration $k$ is:
\begin{equation}
  f_{j}^{(k+1)} = f_{j}^{(k)} \cdot
    \frac{\displaystyle\sum_{i=1}^{I}
      \hat{K}_{i,j}
      \left( n_{\mathrm{meas},i} \,/\, n_{\mathrm{calc},i}^{(k)} \right)}
    {\displaystyle\sum_{i=1}^{I} \hat{K}_{i,j}},
  \label{eq:chahine-update}
\end{equation}
where $\hat{K}_{i,j} = K_{i,j} / \max(\mathbf{K})$ is the normalised kernel,
$n_{\mathrm{meas},i}$ is the measured concentration in the $i$-th voltage bin,
and $n_{\mathrm{calc},i}^{(k)} = \sum_{j} K_{i,j} f_{j}^{(k)}$ is the
forward-calculated concentration at iteration $k$.

The initial guess is a uniform distribution:
\begin{equation}
  f_{j}^{(0)} =
    \frac{\sum_{i=1}^{I} n_{\mathrm{meas},i}}
         {m_{\mathrm{max}} - m_{\mathrm{min}}},
  \quad j = 1, \ldots, J.
  \label{eq:initial-guess}
\end{equation}

%% --------------------------------------------------------------------
\subsection{Markowski 1-2-1 smoothing}
\label{sec:smoothing}
%% --------------------------------------------------------------------

After each iteration, a three-point moving average filter (Markowski, 1987) is
applied to suppress high-frequency oscillations caused by noise amplification:
\begin{equation}
  \tilde{f}_{j} =
    \frac{1}{4} f_{j-1} + \frac{1}{2} f_{j} + \frac{1}{4} f_{j+1},
  \quad j = 2, \ldots, J-1.
  \label{eq:markowski-smoothing}
\end{equation}
The endpoints ($j = 1$ and $j = J$) are not smoothed.

%% ====================================================================
\section{Poisson variance and convergence criterion}
\label{sec:convergence}
%% ====================================================================

%% --------------------------------------------------------------------
\subsection{Poisson counting statistics}
\label{sec:poisson}
%% --------------------------------------------------------------------

The CPC measures particle counts $N_{i}$ during the integration time
$t_{\mathrm{meas},i}$ at each voltage bin.  The measured concentration is
$n_{\mathrm{meas},i} = N_{i} / V_{\mathrm{sample},i}$, where
$V_{\mathrm{sample},i} = Q_{\mathrm{CPC}} \cdot t_{\mathrm{meas},i}$ is the
total sampled air volume.  Because particle counting follows Poisson
statistics, the variance of the measured concentration is:
\begin{equation}
  \mathrm{Var}(n_{i}) =
    \frac{n_{\mathrm{calc},i}}{V_{\mathrm{sample},i}}.
  \label{eq:poisson-variance}
\end{equation}

In practice, a floor is applied to prevent division by zero:
\begin{equation}
  \mathrm{Var}(n_{i}) =
    \max\!\left(
      \frac{n_{\mathrm{calc},i}}{V_{\mathrm{sample},i}}, \;
      \frac{1}{V_{\mathrm{sample},i}^{2}}
    \right).
  \label{eq:variance-floor}
\end{equation}

%% --------------------------------------------------------------------
\subsection{Reduced chi-squared criterion}
\label{sec:chi-squared}
%% --------------------------------------------------------------------

Convergence is determined by the reduced chi-squared statistic:
\begin{equation}
  \chi^{2} = \frac{1}{I} \sum_{i=1}^{I}
    \frac{\left(n_{\mathrm{meas},i} - n_{\mathrm{calc},i}\right)^{2}}
         {\mathrm{Var}(n_{i})}.
  \label{eq:chi-squared}
\end{equation}

The iteration terminates when
$\chi^{2} < \chi^{2}_{\mathrm{threshold}}$.  The default threshold is
$\chi^{2}_{\mathrm{threshold}} = 1.0$, corresponding to the expectation that
residuals should be comparable to the measurement noise.  This criterion
prevents overfitting: stopping at $\chi^{2} \approx 1$ ensures that the
solution explains the data to within the precision allowed by Poisson counting
noise.

%% ====================================================================
\section{Data preprocessing: voltage binning}
\label{sec:preprocessing}
%% ====================================================================

%% --------------------------------------------------------------------
\subsection{Merging upward and downward scans}
\label{sec:scan-merge}
%% --------------------------------------------------------------------

The raw APM data consists of time-series measurements during continuous
voltage scanning (both upward and downward sweeps).  The data are binned into
$I$ voltage bins by grouping measurements with similar applied voltages.

%% --------------------------------------------------------------------
\subsection{Binning rules for physical quantities}
\label{sec:binning-rules}
%% --------------------------------------------------------------------

The binning procedure strictly distinguishes between intensive and extensive
variables:
\begin{itemize}
  \item \textbf{Intensive variables} (state quantities): Applied voltage $V$
    and particle concentration $n$ are averaged within each bin using the
    arithmetic mean.
  \item \textbf{Extensive variables} (amount quantities): The measurement time
    $t_{\mathrm{meas}}$ for each bin is obtained by summing the individual
    time intervals $\Delta t$ of each data row assigned to that bin.
\end{itemize}

This distinction is critical for the correct evaluation of
$V_{\mathrm{sample},i} = Q_{\mathrm{CPC}} \cdot t_{\mathrm{meas},i}$, which
directly enters the Poisson variance calculation~\eqref{eq:variance-floor}.

%% ====================================================================
\section{APM set-point planning with the standalone transfer function}
\label{sec:setpoint}
%% ====================================================================

The RK4 trajectory simulation of Sec.~\ref{sec:apm-transfer} can also be
used on its own, without measurement data, to inspect the APM transfer
function and to plan the operating conditions of an experiment. The tools
described in this section do not modify the inversion pipeline.

%% --------------------------------------------------------------------
\subsection{Standalone transfer function}
\label{sec:standalone-tf}
%% --------------------------------------------------------------------

For a given rotation speed $N$ (in rpm; $\omega = 2\pi N/60$), mobility
diameter $D_{\mathrm{mob}}$ and voltage $V$, the transfer function
$\Omega_{\mathrm{APM}}(m; V)$ is evaluated on a mass grid with exactly the
same code and coefficients as the 1D kernel, so that
$\Omega_{\mathrm{APM}}(m_{j}; V_{i}) = K_{i,j}/\Delta m$. The nominal
classifying mass, at which centrifugal and electrostatic forces balance at
$r_{\mathrm{c}}$ for a singly charged particle, is
\begin{equation}
  m_{\mathrm{c}} = \frac{eV}{\omega^{2} r_{\mathrm{c}}^{2} \ln(r_{2}/r_{1})}.
  \label{eq:classifying-mass}
\end{equation}

%% --------------------------------------------------------------------
\subsection{Resolution parameter and rotation speed}
\label{sec:lambda}
%% --------------------------------------------------------------------

For a spherical particle of diameter $d$ ($= D_{\mathrm{mob}}$) and effective
density $\rho_{\mathrm{eff}}$, the mass, mechanical mobility
$B = Z_{\mathrm{p}}/e$ and relaxation time are
\begin{equation}
  m_{\mathrm{p}} = \frac{\pi}{6} \rho_{\mathrm{eff}} d^{3}, \qquad
  B = \frac{C_{\mathrm{c}}(d)}{3\pi\eta d}, \qquad
  \tau = m_{\mathrm{p}} B.
  \label{eq:sphere-properties}
\end{equation}
With the mean axial velocity
$\bar{v} = Q/[\pi(r_{2}^{2} - r_{1}^{2})]$, the resolution parameter of
Ehara et al.\ (1996) and the resulting angular velocity are
\begin{equation}
  \lambda = \frac{2\tau\omega^{2}L}{\bar{v}}
  \quad\Longrightarrow\quad
  \omega = \sqrt{\frac{\lambda\bar{v}}{2\tau L}}.
  \label{eq:lambda}
\end{equation}

%% --------------------------------------------------------------------
\subsection{Voltage matching the transfer-function mode}
\label{sec:mode-matching}
%% --------------------------------------------------------------------

Equation~\eqref{eq:motion-rearranged} depends on $m$ and $V$ only through
$m\omega^{2}r - eV/[r\ln(r_{2}/r_{1})]$. At fixed $\omega$,
$Z_{\mathrm{p}}$ and $Q$, the transfer function is therefore a function of
$m/V$ alone, and its mode scales linearly with $V$. Starting from the
force-balance voltage
\begin{equation}
  V_{0} = \frac{m_{\mathrm{p}}\omega^{2} r_{\mathrm{c}}^{2}\ln(r_{2}/r_{1})}{e},
  \label{eq:force-balance-voltage}
\end{equation}
the mode $m_{0}^{*}$ of $\Omega_{\mathrm{APM}}(m; V_{0})$ is located
numerically and the voltage is set to
\begin{equation}
  V = V_{0}\,\frac{m_{\mathrm{p}}}{m_{0}^{*}},
  \label{eq:mode-voltage}
\end{equation}
which places the mode exactly at $m_{\mathrm{p}}$. Because the transfer
function has a flat top, the mode is defined as the midpoint of the
contiguous region where
$\Omega_{\mathrm{APM}} \geq 0.99\max\Omega_{\mathrm{APM}}$. The transfer
function is then re-simulated at $(N, V)$ for verification. In practice
$V$ differs from $V_{0}$ by less than $10^{-4}$ (relative) for
$\lambda = 0.2$ and by about $0.5\%$ for $\lambda = 0.5$.

Substituting $V \propto m_{\mathrm{p}}\omega^{2}$ into
\eqref{eq:motion-rearranged}, the right-hand side becomes proportional to
$\tau\omega^{2}/Q \propto \lambda$ times a function of $m/m_{\mathrm{p}}$ and
$r$. Hence, for a given APM geometry, $\Omega_{\mathrm{APM}}$ as a function
of $m/m_{\mathrm{p}}$ depends only on $\lambda$; e.g.\ the peak transmission
and resolution are $0.847$ and $m_{\mathrm{p}}/\mathrm{FWHM} = 2.26$ for
$\lambda = 0.2$, and $0.676$ and $5.15$ for $\lambda = 0.5$ (Kanomax APM-3601
nominal geometry), independent of $d$, $\rho_{\mathrm{eff}}$ and $Q$. The
set-points scale as $N \propto \sqrt{\lambda Q/\tau}$ and
$V \propto m_{\mathrm{p}}\omega^{2}$.

%% --------------------------------------------------------------------
\subsection{Multiply charged sphere clusters}
\label{sec:cluster}
%% --------------------------------------------------------------------

Without a DMA upstream of the APM, a randomly oriented cluster of $n$
primary spheres carrying $q$ elementary charges with $n/q = 1$ has the same
mass-to-charge ratio as the singly charged monomer and is transmitted at the
same set-point. Its mobility is approximated by
\begin{equation}
  d_{\mathrm{ve}} = n^{1/3} d, \qquad
  d_{\mathrm{m}} = \chi\, d_{\mathrm{ve}}, \qquad
  B_{\mathrm{c}} = \frac{C_{\mathrm{c}}(d_{\mathrm{m}})}{3\pi\eta d_{\mathrm{m}}},
  \label{eq:cluster-mobility}
\end{equation}
with the orientation-averaged dynamic shape factor $\chi$ ($1.12$ for a
two-sphere chain; Hinds, 1999, Table~3.2; a continuum-regime value). For
mass $M = n m_{\mathrm{p}}$ and charge $qe$,
\begin{equation}
  \frac{\dd r}{\dd z} \propto q B_{\mathrm{c}}
  \left( \frac{M}{q}\omega^{2}r - \frac{eV}{r\ln(r_{2}/r_{1})} \right),
  \label{eq:cluster-motion}
\end{equation}
so the cluster transfer function is computed with the monomer code using the
voltage $qV$ and the mobility diameter $d_{\mathrm{m}}$. As a function of
$M/q$ it has the shape of a monomer transfer function with
\begin{equation}
  \lambda_{\mathrm{c}} = \frac{2 M B_{\mathrm{c}} \omega^{2} L}{\bar{v}}.
  \label{eq:cluster-lambda}
\end{equation}
For $n = q = 2$, $\lambda_{\mathrm{c}} \approx 1.26$--$1.31\,\lambda$ in the
range $d = 300$--\SI{450}{\nano\meter}, so the cluster transfer function is
narrower and lower than that of the monomer. For monodisperse primaries the
relative counting efficiency at a fixed set-point is
$\Omega_{\mathrm{c}}/\Omega_{1}$ evaluated at $m/q = m_{\mathrm{p}}$; the
ratio of areas $\int\Omega_{\mathrm{APM}}\,\dd(m/q)$ approximates the ratio of
voltage-scan-integrated signals. The actual cluster contribution further
depends on the cluster concentration and charge distribution.

%% --------------------------------------------------------------------
\subsection{Usage}
\label{sec:setpoint-usage}
%% --------------------------------------------------------------------

All settings are given in \texttt{params.py}:
\begin{itemize}
  \item \texttt{TF\_*} (standalone transfer function): rotation speed,
    $D_{\mathrm{mob}}$, list of voltages, and the mass range relative to
    $m_{\mathrm{c}}$. Run \texttt{python run\_transfer\_function.py}.
  \item \texttt{SP\_*} (set-points): list of $(d, \rho_{\mathrm{eff}})$ pairs
    (\si{\nano\meter}, \si{\kilo\gram\per\meter\cubed}), $Q$ (L/min),
    $\lambda$ (a number or a list), mass range relative to
    $m_{\mathrm{p}}$, and output directory. \texttt{SP\_cluster} gives
    $(n, q, \chi)$ of the cluster, or \texttt{None} to skip it. Run
    \texttt{python run\_apm\_setpoint.py}, which writes the transfer
    functions (CSV), the summary tables \texttt{apm\_setpoint\_summary.csv}
    and \texttt{apm\_cluster\_summary.csv}, and the figures (for a list of
    $\lambda$, one figure per particle with all $\lambda$ overlaid).
  \item \texttt{python make\_setpoint\_report.py} then assembles the
    conditions, method, tables and figures into a LaTeX/PDF report
    (\texttt{apm\_setpoint\_report.pdf}) in the output directory.
\end{itemize}
One set-point requires two transfer-function simulations of 401 mass points
(about \SI{40}{\second} on a laptop CPU with $N_{r_{0}} = 1000$ and
$\Delta z = \SI{0.1}{\milli\meter}$); a cluster requires one more.

%% ====================================================================
\section{Implementation mapping}
\label{sec:implementation}
%% ====================================================================

Table~\ref{tab:implementation} maps each theoretical element to the
corresponding function and source file.

\begin{table}[htbp]
  \caption{Mapping from theory to implementation.}
  \label{tab:implementation}
  \centering
  \small
  \begin{tabular}{lll}
    \toprule
    Theory & Implementation & File \\
    \midrule
    Cunningham correction, \eqref{eq:cunningham}
      & \texttt{\_cunningham(Dmob)}
      & \texttt{kernel\_simulator.py} \\
    Trajectory equation, \eqref{eq:motion-rearranged}
      & \texttt{dr\_dz(rad)} inside \texttt{\_rk4\_transmission}
      & \texttt{kernel\_simulator.py} \\
    RK4 integration, \eqref{eq:rk4-k1}--\eqref{eq:rk4-update}
      & \texttt{\_rk4\_transmission(\ldots)}
      & \texttt{kernel\_simulator.py} \\
    Transmission efficiency, \eqref{eq:transmission}--\eqref{eq:flux-weight}
      & \texttt{\_rk4\_transmission(\ldots)}
      & \texttt{kernel\_simulator.py} \\
    1D kernel matrix, \eqref{eq:discrete-forward}
      & \texttt{build\_kernel\_1d(data, params)}
      & \texttt{kernel\_simulator.py} \\
    2D effective kernel, \eqref{eq:2d-kernel}
      & \texttt{build\_kernel\_2d(data, params)}
      & \texttt{kernel\_simulator.py} \\
    Chahine--Twomey update, \eqref{eq:chahine-update}
      & \texttt{solve\_chahine\_twomey(\ldots)}
      & \texttt{inversion\_solver.py} \\
    Markowski smoothing, \eqref{eq:markowski-smoothing}
      & \texttt{solve\_chahine\_twomey(\ldots)}
      & \texttt{inversion\_solver.py} \\
    Poisson variance, \eqref{eq:variance-floor}
      & \texttt{solve\_chahine\_twomey(\ldots)}
      & \texttt{inversion\_solver.py} \\
    Chi-squared criterion, \eqref{eq:chi-squared}
      & \texttt{solve\_chahine\_twomey(\ldots)}
      & \texttt{inversion\_solver.py} \\
    Voltage binning, Sec.~\ref{sec:preprocessing}
      & \texttt{load\_and\_bin(params)}
      & \texttt{data\_parser.py} \\
    Gaussian mode fitting
      & \texttt{fit\_gaussian\_mode(m\_array, f)}
      & \texttt{visualization.py} \\
    Standalone transfer function, \eqref{eq:classifying-mass}
      & \texttt{compute\_apm\_transfer\_function(\ldots)}
      & \texttt{kernel\_simulator.py} \\
    Set-point, \eqref{eq:lambda}--\eqref{eq:mode-voltage}
      & \texttt{find\_apm\_setpoint(\ldots)}
      & \texttt{apm\_setpoint.py} \\
    Cluster, \eqref{eq:cluster-mobility}--\eqref{eq:cluster-lambda}
      & \texttt{simulate\_cluster(\ldots)}
      & \texttt{apm\_setpoint.py} \\
    \bottomrule
  \end{tabular}
\end{table}

%% ====================================================================
\section{References}
\label{sec:references}
%% ====================================================================

\begin{enumerate}
  \item K.~Ehara, C.~Hagwood, and K.~J.~Coakley,
    ``Novel method to classify aerosol particles according to their
    mass-to-charge ratio---Aerosol particle mass analyser,''
    \textit{J.\ Aerosol Sci.}, vol.~27, no.~2, pp.~217--234, 1996.
    DOI: \href{https://doi.org/10.1016/0021-8502(96)00014-4}%
    {10.1016/0021-8502(96)00014-4}.

  \item S.~Twomey,
    ``Comparison of constrained linear inversion and an iterative nonlinear
    algorithm applied to the indirect estimation of particle size
    distributions,''
    \textit{J.\ Comput.\ Phys.}, vol.~18, no.~2, pp.~188--200, 1975.
    DOI: \href{https://doi.org/10.1016/0021-9991(75)90028-5}%
    {10.1016/0021-9991(75)90028-5}.

  \item G.~R.~Markowski,
    ``Improving Twomey's algorithm for inversion of aerosol measurement
    data,''
    \textit{Aerosol Sci.\ Technol.}, vol.~7, no.~2, pp.~127--141, 1987.
    DOI: \href{https://doi.org/10.1080/02786828708959153}%
    {10.1080/02786828708959153}.

  \item W.~C.~Hinds,
    \textit{Aerosol Technology: Properties, Behavior, and Measurement of
    Airborne Particles}, 2nd~ed.\ Wiley, 1999.  Eq.~(3.22) and Table~3.2.
\end{enumerate}

\bigskip
\noindent\textbf{Acknowledgment.}\quad
This document was prepared with the assistance of Claude (Anthropic).
The author assumes full responsibility for the content.

\end{document}
